\FloatBarrier
\section{本章小结}
\label{sec:solution-error-chapter-summary}

本章讨论一个求解结果在多大程度上可信。结果与真解之间的差距有三种来源：问题本身的敏感性、算法在有限精度下引入的误差，以及迭代提前停止留下的误差。本章把它们分开处理，再在问题太敏感时主动改变求解目标，最后在同一组标准下比较前几章的方法。

第一种来源与算法无关。可逆方程的系数矩阵与右端分别有相对扰动$\rho_A$、$\rho_b$，且$\kappa(\bm A)\rho_A<1$时，解的相对误差满足
\[
 \frac{\lVert\Delta\bm x\rVert}{\lVert\bm x_\star\rVert}\leq\frac{\kappa(\bm A)\,(\rho_A+\rho_b)}{1-\kappa(\bm A)\rho_A}.
\]
分子是条件数乘相对扰动，分母是到奇异边界的余量；误差沿小奇异值的方向放大得最多。

最小二乘的三类输出敏感程度不同。拟合值与残差的变化不超过观测的变化，系数却按$1/\sigma_r$放大，所以拟合得好不等于系数可靠。伪逆只在正奇异值有共同下界时连续，在秩增加处不连续，因此数值秩要按误差水平确定，而不是按奇异值是否精确为零确定。

第二种来源由后向误差刻画。残差小只说明后向误差小，前向误差约不超过条件数乘后向误差。后向稳定的算法不能改善病态的问题，但能避免额外的损失：形成正规方程使条件数平方，而QR、SVD与LSQR、LSMR直接作用于$\bm A$。第三种来源要靠核验控制：用真实残差和相应的正规残差判断，并分清达到容差与只是停止计算。

问题本身太敏感时，可以改变求解目标。Tikhonov正则化等价于一个增广最小二乘问题，解唯一当且仅当
\[
 \ker\bm A\cap\ker\bm\Gamma=\{\bm0\}.
\]
岭回归在每个奇异方向上把增益$1/\sigma_i$换成有界的$\sigma_i/(\sigma_i^2+\lambda)$。谱截断只在保留的奇异方向上取倒数，是限制子空间中的最小二乘解，误差分为保留方向上的噪声放大与舍弃方向上的偏差。两者都以拟合误差和偏差换取稳定，参数要联系噪声规模与任务来确定。它们解的是新的目标，较小的系数不能当作把原问题解得更准。

比较方法时，先固定求解目标与验收标准：相容方程看原残差，不相容最小二乘还要看正规残差，最小范数解还要排除零空间分量。再按矩阵的形状、对称正定性、秩与可用的访问方式确定可用的方法。在可用的方法之间，按达到相同精度所需的准备、单步、步数、存储与通信的总成本比较，并把问题病态、算法误差与提前停止三种误差分开。

没有一种方法在所有问题上都更好。直接分解适合能存储因子或有多个右端的问题；Krylov方法适合大型稀疏方程，并按对称正定、对称不定与一般非对称分工；LSQR与LSMR保留矩形最小二乘的结构；行投影法适合只能逐行访问的数据。
